\chapter{Verification: the External Oracle}\label{ch:verify}

\begin{intu}\Intu
Coherence governs whether to keep exploring or stop and reconcile. Only
an external checker touches \emph{truth}. And ``I could not check'' must
never be mistaken for ``it is false''.
\end{intu}

\begin{definition}[Three-valued verification]\label{def:verify}
\textsc{VerifyOp} calls a sandboxed symbolic oracle returning one of
three verdicts, acting on an obstruction counter $\Omega$:
\begin{center}
\renewcommand{\arraystretch}{1.25}
\begin{tabular}{@{}lll@{}}
\toprule
exit & verdict & effect on state\\
\midrule
$0$ & \textsc{Verified} & obstruction cleared; state may crystallise\\
$1$ & \textsc{Refuted} & $\Omega\mathrel{+}=1000$; forced into \textsc{Resolve}\\
$2$ & \textsc{Unverifiable} & $\Omega$ \emph{untouched}\\
\bottomrule
\end{tabular}
\end{center}
\end{definition}

\begin{remark}[Absence of proof is not proof of absence]\label{rem:threevalued}
An earlier version treated every nonzero exit as failure, conflating
\textsc{Refuted} with \textsc{Unverifiable}. That would spike conflict
over the oracle's own blind spots and lock the system into
\textsc{Resolve} over perfectly good mathematics. A checker crash also
returns \textsc{Unverifiable}. This three-valued logic is the promotion
gate of Construction~\ref{con:micro}: a composite becomes a generator of
$\Dop$ only if it survives the oracle, and surviving means
\textsc{Verified}, not merely ``not refuted''.
\end{remark}

\begin{hon}\Hon
The asymmetry in the table is deliberate and is the load-bearing design
decision. A refutation moves $\Omega$ by a large fixed amount; an
inability to check moves it not at all. The system is therefore allowed
to be ignorant without being punished for it, and is punished
immediately for being wrong. The alternative --- treating unverifiable
as mildly bad --- produces a gradient toward only ever asserting things
the oracle happens to understand, which is a narrower system, not a
better one.
\end{hon}

\chapter{The Controller}\label{ch:controller}

\begin{definition}[Coherence-gated mode]\label{def:controller}
With threshold $\varepsilon$, the controller selects
\begin{equation}\label{eq:modes}
  \text{mode}(x)=
  \begin{cases}
    \textsc{Resolve}, & \rhoscore>\varepsilon,\\[2pt]
    \textsc{Explore}, & \rhoscore\le\varepsilon,\\[2pt]
    \textsc{Unknown}, & \rhoscore=\unk .
  \end{cases}
\end{equation}
\end{definition}

\begin{remark}[Three modes, and an empirical $\varepsilon$]\label{rem:controller}
The \textsc{Unknown} branch is essential: when $\rhoscore$ is undefined
the controller \emph{refuses} to default to \textsc{Explore}, because
``no measurement'' is not ``all is well''
(Remark~\ref{rem:unk}). The threshold is calibrated rather than guessed:
on real $384$-dimensional embeddings, agreement gave
$\rhoscore\approx0.00$, related-but-distinct $\approx0.06$, and
off-topic $\approx0.14$, so $\varepsilon\approx0.10$ is a principled
starting gate.

Two cautions were recorded from live runs. First, $\rhoscore$ is
\emph{not} monotonic in the number of active organs: adding a
well-aligned organ raises the normaliser faster than the discord, so
$\rhoscore$ dilutes. Second, comparisons across differently sized
complexes must be made with care, because $\AM$ and $\norm{X}_F^2$ both
change.
\end{remark}

\begin{remark}[Superseded: gate on $\rhotil$, not $\rhoscore$]\label{rem:gate-tilde}
Both cautions are consequences of a single fact, isolated in
Proposition~\ref{prop:split}: $\rhoscore=\dissent\cdot\rhotil$, and
\emph{only $\dissent$ has the pathologies}. Organ-count dilution moves
$\dissent$ and leaves $\rhotil$ invariant. Cross-complex comparison is
unsafe for $\rhoscore$, but $\rhotil$ is comparable across complexes
once expressed as its position inside the window \eqref{eq:window}. The
controller should therefore gate on $\rhotil$ --- or on its window
position --- and report $\dissent$ alongside as the answer to a
genuinely different question.

The typed reading is what the single number could not give:
\begin{itemize}
  \item large $\dissent$, small $\rhotil$: a clean two-camp split; the
        coalition wants restructuring ($\opsplit$, or a new $\opbind$);
  \item small $\dissent$, large $\rhotil$: one organ is an outlier; a
        single edge wants repair.
\end{itemize}
\end{remark}

\begin{workedbox}\Worked
Apply the controller to the running example. With $\varepsilon=0.10$ and
$\rhoscore=0.2$, mode \eqref{eq:modes} returns \textsc{Resolve}, and a
diagnostic reading only $\rhoscore$ would report ``mild disagreement,
reconcile''. The typed reading gives more: $\dissent=0.5$ is large and
$\rhotil$ sits at window position $\tfrac13$, so by
Remark~\ref{rem:gate-tilde} this is a two-camp split, and the correct
remediation is $\opsplit$ or a new $\opbind$ --- \emph{not} repair of the
worst edge. Note also that the worst-edge diagnostic was ambiguous here
(the argmax tied), which is the same fact appearing in a second
instrument.
\hfill$\square$
\end{workedbox}

\begin{hon}\Hon
The controller is the one place where every earlier honesty decision
becomes operational. $\unk$ propagates into a mode rather than being
coerced to a number; the split of Proposition~\ref{prop:split} changes
which quantity is gated on; the window of
Proposition~\ref{prop:window} makes the threshold topology-relative. It
is also where a mistake is most expensive, because a controller that
mis-reads its instrument does not merely report wrongly --- it acts.
\end{hon}
